\chapter{绪论}

\section{研究背景与意义}

证券担保品横跨权益、固定收益、基金、衍生品和信用风险合约。六类产品的报价机制、价值口径、现金流结构、交易日历和到期事件差异显著：股票价值受公司行为、停牌和涨跌停影响；债券需要同时处理全价、应计利息、票息、本金和违约；基金可能以净值或可交易价格计量并受到申赎和折溢价影响；期货具有合约乘数、逐日盯市、保证金和展期规则；期权价值由标的价格、期限、波动率曲面和行权事件共同决定；信用衍生品还涉及保护费、信用事件、回收率和结算方式。将这些产品简单压缩为同一种价格预测，既会丢失产品经济含义，也会使风险结果无法与真实损益一致。

担保品管理关注的并非只有“今天值多少”，还包括在预定持有或处置期内价值如何演化。设估值时点统一价值为 \(V_t\)，目标时点价值为 \(V_{t+h}\)，期内现金流为 \(C_{t,t+h}\)，则价值差与总损益分别为
\begin{align}
  \Delta V_{t,h} &= V_{t+h}-V_t,\\
  \mathrm{PnL}_{t,h} &= V_{t+h}-V_t+C_{t,t+h}.
\end{align}
与损益口径一致的正损失定义为 \(L_{t,h}=-\mathrm{PnL}_{t,h}\)；当明确采用不含现金流的价值口径时，才使用 \(L^{\mathrm{value}}_{t,h}=V_t-V_{t+h}\)。模型需要预测 \(V_{t+h}\) 的点值和分位数以及 \(L_{t,h}\) 的完整或近似分布，由此导出 VaR、ES、预测区间和压力损失。订单簿买卖价差、信用利差、基金折溢价、期货基差和期权波动率偏斜仍有解释价值，但只作为产品特有特征或辅助诊断。

这一口径使六类产品可以在统一问题上比较，同时保留各自的定价约束和现金流机制。科学上，它要求模型从单点拟合转向条件价值路径和损失分布预测，并使用 proper scoring rule、波动率损失、覆盖检验和独立性检验评价不确定性；工程上，它要求每一次预测都绑定真实的 \texttt{as\_of} 信息集、产品条款、数据版本和模型版本；验收上，它使 VaR 与随后实现的 PnL 能够在同一损失基础上进行返回检验，避免用方向准确率、\(1-\mathrm{MAPE}\) 或某个局部指标替代合同风险结论。

多源数据和语义信息为估值及风险预测提供了新的条件变量，但也引入发布时间错配、数据修订、文本泄漏和不可解释输出等风险。本课题因此不把大模型作为脱离数值底座的价格生成器，而将其限定为有边界、有对照、有回退的残差或状态增强器。只有真实时点语义在相同样本上稳定优于无文本和打乱/错位文本对照，且不显著恶化合同风险指标，才认为语义信息产生了增量技能。

本研究的意义体现在三个方面。其一，形成跨产品一致的“当前价值—未来价值—持有期损益—尾部风险”研究对象，使估值、价差风险和返回检验不再相互脱节。其二，为六类担保品建立产品级主模型、共同输出协议和可比较的基准阶梯，使模型数量、指标数量与算法组件各归其位。其三，将数据时点、实验预注册、模型回退、异常诊断和证据导出纳入测试系统，使研究结论能够被重跑、核对和审计，并为后续第三方测试和机构应用提供可定位证据。

\section{研究需求与目标}

本课题以《项目申报指南》《项目/课题任务书》《最终考核指标说明》及已经确认的项目边界为共同约束；当原始文件表述不完全一致时执行要求并集。总体目标是形成一套覆盖六类证券担保品、同时输出估值与持有期价差风险、能够进行至少五年实际市场返回检验并支持当前数据监控的可审计技术体系。具体目标如下。

\begin{enumerate}[label=\arabic*)]
  \item \textbf{建立统一计量对象与输出契约。} 对每个资产和预测时点显式记录产品类型、产品标识、\texttt{as\_of} 日期、目标日期、持有期、价值口径、期内现金流、损失基础和质量标志；统一输出估值点值及分位数、未来价值、价值差、PnL、损失分布、波动率、VaR、ES、流动性指标和产品敏感度。暂定同时报告 T+1、T+5、T+10，主持有期在书面确认前不从结果中选择。
  \item \textbf{形成六个产品级估值及价差风险主模型。} M1--M6 分别覆盖全部 A 股、公募债券、公募基金、场内期货、场内期权和以场内证券为标的的信用衍生品。每一类均配置朴素基准、产品经典理论基准和强现代基准；只有在目标、样本和信息集可比时，才加入代表性近期方法作对照。六个主模型均需给出至少六项风险或流动性输出，其中包括波动率和 VaR。
  \item \textbf{实现两个受约束的大模型增强器。} L1（TFHTS 估值残差增强器）检验时点可得文本是否改善六类底座的估值与收益分布；L2（EPU--PanelMLP 持有期风险增强器）检验政策不确定性及其他语义状态是否改善未来价值、损失分位数、波动率和 VaR。两个增强器均需报告输出覆盖率、实际增强率、通过率和回退率，不能把底座回退静默计为大模型增量成功。
  \item \textbf{建设实际市场结果验证与动态监控系统。} 系统应支持冻结数据快照、训练与预测编排、至少五年历史回放、VaR/PnL 返回检验、当前数据预测、漂移与异常监控、统一接口、权限与审计日志，并能够向第三方测试和应用示范导出不可手改的证据包。系统还需验证安全、稳定、并发、异常恢复、参数化配置、定制指标与报表以及外部课题接口。
  \item \textbf{建立分层验收与证据追踪。} 对合同验收、科学技能、工程质量和产品覆盖分别给出状态。合同“准确率不低于 \(90\%\)”以五年窗口下 VaR 与实际 PnL 的返回检验为主线，但其置信水平、容差、资产聚合和样本通过公式在签字前只输出暂定结果；估值点预测使用 MAE、RMSE、WAPE 或 sMAPE、MASE 等指标作科学评价，不用单一 \(1-\mathrm{MAPE}\) 代替合同结论。
\end{enumerate}

\section{面临的关键问题}

\textbf{（1）统一研究对象与产品异质性之间的协调。}
六类产品必须回答同一个持有期价值与损失问题，但不能采用同一种粗糙的价格差定义。研究需要在统一输出契约下分别处理复权、票息与应计、净值与市场价、逐日盯市与展期、波动率曲面与行权、信用事件与回收等机制，并保证预测损失与实现 PnL 使用相同口径。

\textbf{（2）持有期价差风险及合同返回检验的精确定义。}
现有原始文件明确要求比较 VaR 与实际损益并主要采用五年历史窗口，却没有唯一给出主持有期、VaR 置信水平、“\(90\%\)”的统计含义、合理置信区间和六类聚合规则。对于 \(h>1\)，相邻预测的实现窗口重叠，若直接对全序列进行独立性检验会产生机械相关。因此需要事前冻结损失正负号、现金流基础、T+1/T+5/T+10 的主次关系，并按 \(t\bmod h\) 构造非重叠 cohort 进行独立性检验。

\textbf{（3）时点可得信息和数据泄漏控制。}
市场、基本面、宏观、政策和文本数据的记录日期不等于真实可用日期；修订数据、回填数据库、随机切分、全样本标准化和重叠目标窗口均可能产生虚高结果。所有特征必须保留来源时间和最早可用时间，只允许按时间扩展或滚动切分，并在目标窗口重叠时实施 purge/embargo。缺失填补、标准化、词表、特征选择和超参数选择只能在训练期拟合。

\textbf{（4）全品种覆盖分母和信用衍生品数据缺口。}
“全部”不能由抽样实验或运行成功的子集代表。六类产品需要逐日生成 master catalogue、eligible catalogue、预测台账和排除台账，分别报告数据资格率、运行覆盖率和可回测率。当前以场内证券为标的的信用衍生品缺少真实端到端合约、报价、成交、回收率和信用事件数据；在数据或书面替代决定到位前，M6 只能形成阻塞证据，合成数据只能验证工程链路。

\textbf{（5）科学增益与合同通过的区分。}
经验覆盖率达到某一阈值不等于损失分布预测优于合理基准，估值误差较小也不自动代表 VaR 校准正确。每个产品必须在相同 eligible sample 和时点信息集上比较朴素、经典和强现代基准，并同时报告绝对误差、相对差异、置信区间、压力期结果和失败案例。合同通过、科学通过、工程通过和覆盖通过必须分别记录。

\textbf{（6）大模型语义增量的可识别性与安全回退。}
文本模型可能只学习时间、市场状态或数据处理差异，而非真正利用语义。L1/L2 必须设置 no-text、shuffled 和 time-shifted 对照，保存检索证据、模型版本和调用记录；文本缺失、延迟、冲突、提示注入或输出越界时，系统应确定性回退到产品底座。未证明增量的实现可保留，但状态只能是“方法已实现、未证明增量”。

\textbf{（7）研究结果、系统输出与验收材料的一致性。}
手工复制表格、事后调整分母或用平台截图替代原始结果，会再次造成模型、报告和测试大纲分叉。研究链路必须由冻结配置生成预测明细、返回检验、模型卡、数据卡、失败清单和哈希 manifest，报告中的主要数字只能从这些证据反向生成。第三方测试、三家机构应用、接口联调和安全稳定性属于外部或工程验收，不得由模型实验结果代替。

\section{本文研究内容与逻辑关系}

\subsection{主要研究内容与关键技术}

本课题的主要研究内容按成果对象组织为三条研究主线，并由系统验证和应用验证闭合证据链。

\textbf{研究点一：多品种证券担保品估值及价差风险计量技术。}
围绕六类产品建立 M1--M6 主模型。M1 以可识别的 GBSV 家族为股票价格与波动率分布研究核心；M2 以现金流、收益率曲线和无套利价值路径为债券主干；M3 以净值、持仓穿透和市场因子价值路径为基金主干；M4 以持有成本、期限结构和结算价路径为期货主干；M5 以 Black--Scholes/Black--76、隐含波动率曲面和 Greeks 为期权主干；M6 以违约强度、生存概率、回收率和合约现金流为信用衍生品主干。各模型由同一条件价值路径生成未来价值、持有期 PnL、损失分布、VaR 和 ES。

\begin{table}[htbp]
  \centering
  \caption{六类产品主模型的统一定位}
  \label{tab:m1-m6-overview}
  \ReportTableText
  \begin{tabularx}{\textwidth}{@{}p{3.2em}p{7.2em}X@{}}
    \toprule
    模型 & 产品 & 估值与持有期风险主线 \\
    \midrule
    M1 & 全部 A 股 & GBSV 条件价格/波动率路径，导出价值差、损失分布、VaR 与 ES \\
    M2 & 公募债券 & 现金流与收益率曲线定价，纳入应计、票息、本金和信用事件 \\
    M3 & 公募基金 & 净值或市场价与持仓/因子路径，纳入分红、申赎、折算和清盘 \\
    M4 & 场内期货 & 持有成本与期限结构路径，纳入乘数、逐日盯市、到期和展期 \\
    M5 & 场内期权 & 标的与波动率曲面联合演化，纳入权利金、行权、到期和 Greeks \\
    M6 & 信用衍生品 & 违约强度、生存/回收和双腿现金流；真实数据到位前保持阻塞 \\
    \bottomrule
  \end{tabularx}
\end{table}

\textbf{研究点二：大模型增强的估值及价差风险计量技术。}
在 M1--M6 底座上构建共享的语义编码、检索、约束、校准和回退框架。L1 输出估值点值或分位数的受限残差及状态权重；L2 输出未来价值、损失分位数、波动率或 VaR 的受限残差/状态门控。关键技术包括真实发布时间对齐、检索证据绑定、参数高效适配、无套利或产品边界约束、消融试验和确定性回退。

\textbf{研究点三：实际市场结果返回检验与动态监控技术。}
围绕五年历史回放和当前 \texttt{as\_of} 预测，建立统一标签、时间切分、VaR/PnL 返回检验、分布与波动率评分、压力期分层、漂移检测和异常诊断方法。关键技术包括重叠持有期的 cohort 检验、日期块 bootstrap/HAC、catalogue 覆盖核算、实验预注册、不可变数据快照和证据 manifest。

\textbf{系统与应用验证。}
成果3将三条研究主线实现为可配置、可编排、可回放、可监控和可审计的测试系统；随后分别开展功能、性能、安全、稳定、接口和异常恢复验证，以及不少于三家机构的应用示范和第三方检测。系统验证回答“能否可靠运行”，应用验证回答“能否在约定场景形成可接受证据”，二者均不替代模型科学评价。

\subsection{研究内容的逻辑关系}

三条研究主线共享同一时点数据契约和预测输出协议。研究点一提供六类产品的受约束数值底座；研究点二只在底座之上产生可检验的语义残差或状态修正；研究点三对底座与增强版本采用相同样本、目标和返回检验，并将结果交给测试系统监控和归档。系统验证与应用验证使用同一预测台账和证据包，避免另造一套验收数字。

\begin{figure}[htbp]
  \centering
  \fbox{%
    \begin{minipage}{0.91\textwidth}
      \centering
      时点可得的市场、宏观、主体、文本与产品条款数据\\[0.6ex]
      \(\Downarrow\)\\[-0.2ex]
      六类产品底座 M1--M6：当前价值、未来价值与持有期损失分布\\[0.6ex]
      \(\Downarrow\)\\[-0.2ex]
      L1/L2 语义增强：受限残差、消融对照、质量门禁与回退\\[0.6ex]
      \(\Downarrow\)\\[-0.2ex]
      统一评价：估值误差、分布评分、波动率、VaR/ES 与五年返回检验\\[0.6ex]
      \(\Downarrow\)\\[-0.2ex]
      测试系统：编排、监控、权限、审计、接口与证据导出\\[0.6ex]
      \(\Downarrow\)\\[-0.2ex]
      第三方测试和机构应用：合同、科学、工程、覆盖四类状态分别判定
    \end{minipage}%
  }
  \caption{本课题研究内容及证据闭环}
  \label{fig:research-logic}
\end{figure}

该关系同时规定了结论的边界：没有真实 catalogue 不能声称全覆盖，没有冻结返回检验公式不能声称合同通过，没有对照消融不能声称大模型增量，没有可重跑 manifest 不能声称工程结果可审计。某一层失败时保留失败证据并阻止更高层结论，但不删除已经验证的下层能力。

\section{本文组织结构}

本报告共设九章。第 0 章为引言，界定研究对象、公共价差风险口径、成果边界和证据原则。第 1 章为绪论，说明研究背景与意义、需求与目标、关键问题、主要研究内容及其逻辑关系。第 2 章综述多品种担保品估值与风险计量、大模型金融时序增强以及返回检验和动态监控技术的研究进展，并据此定位本课题的研究切入点。

第 3 章研究六类证券担保品的产品级估值与持有期损失分布模型，给出统一输出契约、M1--M6 的产品特有机制、基准阶梯和实验结果。第 4 章研究 L1/L2 大模型增强方法，重点报告时间对齐、检索证据、约束与回退机制以及 no-text、shuffled、time-shifted 消融。第 5 章研究实际市场结果返回检验与动态监控方法，给出五年历史回放、重叠持有期检验、覆盖核算、压力期诊断和证据生成规则。

第 6 章在纳入交付范围时报告测试系统的设计、实现及功能、性能、安全、稳定和接口验证。第 7 章在外部条件具备时报告机构应用设计、实施过程、第三方测试和示范结果。第 8 章总结主要研究发现、创新边界、未完成事项和未来研究方向。其后依次列出参考文献，以及研究期间发表的学术论文与其他相关学术成果。
